\documentclass[tikz,border=10pt]{standalone}
\usepackage{tikz}
\usetikzlibrary{calc}

\begin{document}
\begin{tikzpicture}[scale=1.2, line cap=round, line join=round]
    % 1. 定义核心点坐标 (基于图像视觉比例)
    \coordinate (G) at (0,0);           % 地面辅助线左端点
    \coordinate (D) at (3,0);           % 地面辅助线右端点/直角点
    \coordinate (A) at (0,1.2);         % 点A的位置
    \coordinate (B) at (0,5);           % 点B，位于A上方
    \coordinate (C) at (3,0.7);         % 点C，位于D上方
    \coordinate (E) at (2.2,0.6);       % 点E，在斜线上
    
    % 2. 绘制水平地面 (左右略宽出辅助线范围)
    \draw[thick] (-0.8,0) -- (G);
    \draw[thick] (D) -- (4,0);
    
    % 3. 绘制实线部分 (图形主体)
    \draw[thick] (B) -- (A) -- (E) -- (C) -- (D); % 连接 B-A-E-C-D
    \draw[thick] (B) -- (E);                       % 连接 B-E
    \draw[thick] (E) -- (D);                       % 连接 E-D (形成 A-E-D 直线感)
    
    % 4. 绘制虚线部分 (辅助线)
    \draw[dashed, thick] (A) -- (G);               % 垂线 AG
    \draw[dashed] (G) -- (D);               % 地面辅助线段 GD
    
    % 5. 标注字母节点
    % anchor 用于微调文字相对于坐标点的方位
    \node[left]  at (A) {$A$};
    \node[above] at (B) {$B$};
    \node[above right] at (C) {$C$};
    \node[below] at (D) {$D$};
    \node[above right] at (E) {$E$};
    \node[below] at (G) {$G$};
    
    % 6. (可选) 如果需要标注直角符号，可以取消下面代码的注释
    % \draw (G) +(0,0.2) -- +(0.2,0.2) -- +(0.2,0);
    % \draw (D) +(-0.2,0) -- +(-0.2,0.2) -- +(0,0.2);

\end{tikzpicture}
\end{document}